% introduction_constructs.tex
\subsection{Constructs and the Nomological Network}

Psychological constructs cannot be inspected directly; they are postulated attributes inferred from regularities in what people do, feel, and report. \citet{cronbach-meehl-1955} formalized this as construct validity: where no criterion or content universe fully defines an attribute, a test is validated only by the construct's place in a nomological network of lawful relations to indicators, experimental manipulations, and other constructs. Stating a construct's laws is saying what it is, elaborating them is learning more, and a sparse network leaves a vague construct: meaning is relational---the evidential web's breadth and coherence, not a hidden essence. \citet{messick-1995} recast all validation as construct validation---an integrated evaluative judgment about score interpretations, not a property of instruments---with six aspects of evidence (content, substantive, structural, generalizability, external, consequential) and two pervasive failure modes: construct underrepresentation and construct-irrelevant variance. \citet{clark-watson-2019} make the network a working tool: a precise written conceptualization, boundaries with near neighbors included, is scale development's first and most consequential step, for without an articulated theory there is no construct validity to investigate. \citet{borsboom-etal-2004} dissent on ontological grounds: a test is valid only if its attribute exists and causally produces variation in test outcomes, and psychology's networks are too loosely specified to fix the meanings of theoretical terms. We adopt the relational framework but concede the realist point: no correlational pattern, behavioral or semantic, by itself certifies that an attribute exists; both structures figure here as evidence about how a construct is defined and how its indicators behave, not as proof of its reality.
